\documentclass[12pt,a4paper]{article}
\usepackage[UTF8]{ctex}
\usepackage{geometry}
\usepackage{amsmath,amssymb}
\usepackage{graphicx}
\usepackage{booktabs}
\usepackage{longtable}
\usepackage{array}
\usepackage{enumitem}
\usepackage{hyperref}
\usepackage{fancyhdr}
\usepackage{xcolor}
\usepackage{colortbl}
\usepackage{setspace}

\geometry{left=2.5cm,right=2.5cm,top=2.5cm,bottom=2.5cm}
\setlength{\headheight}{15pt}
\setstretch{1.5}

% 去掉页眉中的“现有技术检索报告”
\pagestyle{fancy}
\fancyhf{}
\fancyfoot[C]{\thepage}
\renewcommand{\headrulewidth}{0pt}

\begin{document}

% 封面/头部信息区
\noindent

\vspace{1em}

\begin{center}
    {\Huge\heiti 中国移动专利申请} \\[0.5em]
    {\Huge\heiti 检索报告} \\[1em]
\end{center}

\begin{table}[h]
    \centering
    \renewcommand{\arraystretch}{2}
    \begin{tabular}{|p{3cm}|p{12cm}|}
        \hline
        \textbf{发明名称} & \textcolor{blue}{噪声自适应变分量子加密（NAV-QE）系统、方法及其设备指纹与硬件绑定密钥应用} \\
        \hline
        \textbf{申报单位} & \textcolor{blue}{中国移动通信有限公司研究院} \\
        \hline
        \textbf{检索人} & \textcolor{blue}{【请填写实际检索人姓名/部门】} \\
        \hline
        \textbf{检索日期} & \colorbox{yellow}{2026.03.06} \\
        \hline
        \textbf{关联项目} & （待填写） \\
        \hline
    \end{tabular}
\end{table}

\vspace{0.5em}

\noindent\fbox{\parbox{0.97\textwidth}{
\textbf{与量子计算的关联说明：}

本申请聚焦于将量子处理器在运行和校准过程中呈现的固有噪声特征（如 $T_1$、$T_2$、门误差、串扰、读出误差等）转化为硬件指纹、证明材料与设备绑定密钥。相较于单纯量子随机数生成、传统量子误差表征或经典 PUF，本申请强调“噪声特征提取 + 指纹量化映射 + 密钥导出 + 篡改检测 + 云/控制面部署”的整体技术链条，检索时需同时覆盖量子噪声表征、硬件证明、PUF、量子云认证、后量子密钥派生等交叉方向。
}}

\vspace{1em}

% 正文部分 - 严格按照模板结构

% 二、使用的中文与外文检索关键词
\section*{一、使用的中文与外文检索关键词}

\textbf{中文检索关键词：}
\begin{enumerate}[leftmargin=2em]
    \item 量子噪声；量子硬件指纹；量子处理器指纹；量子设备认证；
    \item 变分量子电路安全；参数化量子电路；噪声表征；退相干特征；
    \item $T_1$、$T_2$、门误差、串扰、读出误差；
    \item 量子物理不可克隆函数；量子 PUF；硬件证明；
    \item 量子云可信执行；量子控制平面；量子中间件；
    \item 后量子密钥派生；设备绑定密钥；篡改检测；Mahalanobis 距离。
\end{enumerate}

\textbf{外文检索关键词：}
\begin{enumerate}[leftmargin=2em]
    \item quantum noise; quantum hardware fingerprinting; quantum processor fingerprint; device attestation;
    \item variational quantum circuit security; parameterized quantum circuit; noise characterization;
    \item T1 relaxation; T2 dephasing; gate error; crosstalk; readout error;
    \item quantum physical unclonable function; quantum PUF; hardware-bound key;
    \item quantum cloud attestation; trusted quantum computing; quantum middleware;
    \item post-quantum key derivation; tamper detection; Mahalanobis distance.
\end{enumerate}

% 三、相关专利文献
\section*{二、相关专利文献}

\renewcommand{\arraystretch}{1.5}
\begin{longtable}{|p{1.5cm}|p{6cm}|p{7.5cm}|}
\hline
\textbf{编号} & \textbf{相关度类别} & \textbf{文献信息 (标题/来源/作者/日期)} \\
\hline
1 & A（基础技术） & \textbf{Scalable and Robust Randomized Benchmarking of Quantum Processes} \newline Physical Review Letters, 2011 \newline Easwar Magesan, Jay M. Gambetta, Joseph Emerson \\
\hline
2 & A（基础技术） & \textbf{Prescription for Experimental Determination of the Dynamics of a Quantum Black Box} \newline Journal of Modern Optics, 1997 \newline Isaac L. Chuang, Michael A. Nielsen \\
\hline
3 & Y（相关技术） & \textbf{Quantum Physical Unclonable Functions} \newline arXiv:1905.02550, 2019 \newline Myrto Arapinis et al. \\
\hline
4 & A（背景技术） & \textbf{Physical One-Way Functions} \newline Science, 2002 \newline Ravikanth Pappu et al. \\
\hline
5 & A（背景技术） & \textbf{Quantum Random Number Generators} \newline Reviews of Modern Physics, 2017 \newline Carlos A. F. Herrero-Collantes, Juan C. Garcia-Escartin \\
\hline
\end{longtable}

% 四、分析评述
\section*{三、分析评述}

\textbf{1. 与文献1、文献2（量子噪声表征基础文献）的对比分析：}
文献1和文献2公开了量子门误差、量子过程及相关噪声参数的测量与拟合方法，为量子处理器表征提供了基础工具，但其目标在于提升量子计算精度、校准质量和误差认知，并未公开将这些噪声参数作为安全资源来构建设备指纹、导出设备绑定密钥或执行持续证明。

\textbf{2. 与文献3（量子 PUF）的对比分析：}
文献3从量子态或量子随机过程的不可克隆性角度讨论量子 PUF 概念，具有一定启发意义，但未明确公开利用实际 NISQ 量子处理器已有校准噪声参数构建可重构指纹，也未公开机器学习表征、量化映射、后量子密钥派生和云/控制面部署架构。

\textbf{3. 与文献4（经典 PUF）的对比分析：}
文献4是经典 PUF 方向的基础工作，其物理基础是经典制造差异，不是量子退相干、门误差和串扰，因此不能直接迁移到量子处理器场景。本申请则利用量子硬件原生噪声这一不同物理来源作为身份根，并进一步处理连续噪声的稳定重构问题。

\textbf{4. 与文献5（量子随机数生成）的对比分析：}
文献5主要解决高熵随机数生成问题，不涉及设备身份绑定、后量子密钥绑定或篡改检测。本申请不仅利用量子源获得熵，还强调设备唯一性、可重复重构性和部署级安全联动。

% 五、检索结论
\section*{四、检索结论}

经检索，暂未发现与本申请全部技术特征相同的现有技术。本申请提出的基于量子处理器固有噪声特征的硬件指纹、设备绑定密钥与篡改检测方案，在“噪声特征提取—指纹量化映射—密钥派生—持续检测—工程部署”这一完整闭环上具有较好的新颖性和创造性基础。

\end{document}
